% SPDX-License-Identifier: MIT
\chapter{Quick start}
\label{ch:quickstart}

This chapter produces a finished, annotated figure in about five minutes. It uses
the example data in the repository, so you can follow it exactly before
substituting your own tree.

\section{Your first figure, from the command line}

The fastest route to a figure needs no interface at all:

\begin{lstlisting}[language=bash]
makeyourtree render examples/primates.nwk figure.svg
\end{lstlisting}

That writes a rectangular phylogram as a vector SVG. Add a layout and an
annotation track:

\begin{lstlisting}[language=bash]
makeyourtree render examples/primates.nwk figure.svg \
    --mode circular --align-tips \
    --track examples/primates_region.mytrack
\end{lstlisting}

\begin{figure}[htbp]\centering
  \includegraphics[width=0.72\linewidth]{layout-circular}
  \caption{A circular fan, the output of the second command above (shown
  without the track for clarity).}\label{fig:quickstart-circular}
\end{figure}

\section{The same thing in the application}

\begin{enumerate}
  \item Start the application:
    \begin{lstlisting}[language=bash]
makeyourtree-studio
    \end{lstlisting}
  \item \textbf{File \textbf{\textrightarrow} Open} and choose
        \cmd{examples/primates.nwk}. The tree appears immediately.
  \item Press \keys{3} to switch to a circular layout, or \keys{1} to go back to
        rectangular. The layout keys are \keys{1}--\keys{5}
        (Appendix~\ref{ch:shortcuts}).
  \item Press \keys{A} to align the tips and draw guide lines.
  \item \textbf{Annotate \textbf{\textrightarrow} Import Annotations}
        (\keys{Ctrl+I}) and choose \cmd{examples/primates\_region.mytrack}. A
        colour strip appears beside the tips, with a legend.
  \item \textbf{File \textbf{\textrightarrow} Export Figure} (\keys{Ctrl+E}).
        Choose SVG, PDF or PNG; the dialog predicts the pixel dimensions before
        you commit.
\end{enumerate}

Every step is undoable with \keys{Ctrl+Z}.

\section{A publication figure}

The example data includes a 24-tip bacterial tree with two annotation files.
Loading the tree and both tracks gives Figure~\ref{fig:combined-rectangular}:

\begin{lstlisting}[language=bash]
makeyourtree render examples/bacteria.nwk figure.svg --align-tips \
    --track examples/bacteria_resistance.mytrack \
    --track examples/bacteria_expression.mytrack
\end{lstlisting}

\galleryfig{combined-rectangular}{Three stacked tracks --- a categorical colour
strip, a diverging heatmap and a bar chart with a value axis --- beside a
rectangular phylogram. Tips are aligned and guide lines connect them to their
labels.}{1.0}

Switching that same figure to a circular layout requires changing one option, and
no track behaves differently:

\galleryfig{combined-circular}{The identical tracks and data in a circular
layout. The heatmap has become an annulus and the bars radiate outward; no track
code differs between this figure and Figure~\ref{fig:combined-rectangular}. This
is the band-space design of Section~\ref{sec:design} doing its work.}{0.85}

\section{From Python}

The same figure, reproducibly, from a script:

\begin{lstlisting}[language=Python]
from makeyourtree.io import load_tree
from makeyourtree.doc import Document
from makeyourtree.layout import LayoutParams, LayoutMode
from makeyourtree.scene import compose
from makeyourtree.render import render_svg
from makeyourtree.ops import midpoint_root, ladderize

tree = load_tree("examples/primates.nwk")
midpoint_root(tree).apply(tree)
ladderize(tree).apply(tree)

doc = Document(tree=tree,
               params=LayoutParams(mode=LayoutMode.CIRCULAR, arc=340))
with open("figure.svg", "w", encoding="utf-8") as fh:
    fh.write(render_svg(compose(doc)))
\end{lstlisting}

Chapter~\ref{ch:python} covers the API in full.

\section{Where to go next}

\begin{itemize}
  \item Chapter~\ref{ch:formats} --- what \app{} can read, and what it does with
        files that break their own specification.
  \item Chapter~\ref{ch:layouts} --- every layout mode, with figures.
  \item Chapter~\ref{ch:annotation} --- writing your own annotation files.
  \item Chapter~\ref{ch:tracks} --- the reference for all thirteen track types.
\end{itemize}
